\documentclass[11pt]{article}

\usepackage{amsmath, amssymb}
\usepackage{geometry}
\usepackage{natbib}
\usepackage{hyperref}
\usepackage{tcolorbox}

\geometry{margin=1in}

\newtcolorbox{intuitionbox}{
  colback=gray!10,
  colframe=black,
  title=Intuition,
  fonttitle=\bfseries
}

\newtcolorbox{technicalbox}{
  colback=white,
  colframe=black,
  title=Technical Statement,
  fonttitle=\bfseries
}

\title{More Structural Shocks than Observables in DSGE Estimation } % \\ \large A Pedagogical Note
\author{}
\date{}

\begin{document}

\maketitle

\section{State-Space Representation}

\begin{intuitionbox}
Structural shocks do not map directly into observables. 
They affect latent state variables, which then propagate dynamically through the model before influencing data. 
Therefore, counting shocks and observables is not a meaningful identification criterion.
\end{intuitionbox}

\begin{technicalbox}
A linearized DSGE model can be written in state-space form:
\begin{align}
x_t &= F(\theta) x_{t-1} + G(\theta) \varepsilon_t, 
\quad \varepsilon_t \sim \mathcal{N}(0,Q), \\
y_t &= H(\theta) x_t + u_t,
\quad u_t \sim \mathcal{N}(0,R).
\end{align}
Here $x_t$ is the vector of latent states and $y_t$ the observables.
Shocks influence observables only indirectly via the transition matrix $F$.
\end{technicalbox}

References: \citet{an2007}, \citet{herbst2015}, \citet{durbin2012}.

\section{The Likelihood and the Kalman Filter}

\begin{intuitionbox}
The Kalman filter computes the likelihood by integrating over unobserved state variables.
You do not estimate shocks directly from observables.
You estimate a dynamic system and marginalize out latent states.
\end{intuitionbox}

\begin{technicalbox}
The likelihood integrates over latent states:
\begin{equation}
p(y_{1:T}\mid\theta)
=
\int p(y_{1:T},x_{1:T}\mid\theta)\, dx_{1:T}.
\end{equation}

Under Gaussian assumptions, the log-likelihood has the prediction-error form:
\begin{equation}
\log L(\theta)
=
\sum_{t=1}^{T}
\left[
-\frac{1}{2}\log|S_t|
-\frac{1}{2}\nu_t' S_t^{-1}\nu_t
\right]
+ \text{constant},
\end{equation}
where $\nu_t$ is the one-step-ahead forecast error and $S_t$ its variance.
\end{technicalbox}

Key references: \citet{durbin2012}, \citet{hamilton1994}.

\section{Why Medium-Scale NK Models Can Sustain Many Shocks}

\begin{intuitionbox}
Identification comes from dynamic structure, not counting dimensions.
If two shocks generate different impulse responses, the data can distinguish them.
\end{intuitionbox}

\begin{technicalbox}
Identification requires that the mapping from parameters and shock processes
into the likelihood-implied covariance structure be locally invertible.
Distinct persistence parameters and distinct propagation through $F$ generate
linearly independent impulse responses.
\end{technicalbox}

Reference: \citet{smets2007}.

\section{When ``Too Many Shocks'' Becomes a Problem}

\begin{intuitionbox}
Problems arise when shocks look the same in the data or absorb misspecification.
If two shocks generate nearly identical dynamics, identification weakens.
\end{intuitionbox}

\begin{technicalbox}
Weak identification emerges when the Jacobian of the mapping from parameters
to likelihood moments becomes rank deficient.
Variance swapping and prior sensitivity may occur.
See \citet{iskrev2010}.
Misspecification concerns are emphasized in \citet{denhaan2021asd}.
Correlated disturbances are studied in \citet{curdia2012}.
\end{technicalbox}

\section{Identification Criterion}

\begin{intuitionbox}
The relevant question is not:
``Are there more shocks than observables?''
The relevant question is:
``Are the shocks dynamically distinguishable?''
\end{intuitionbox}

\begin{technicalbox}
The relevant condition is not:
\begin{equation}
\#\text{shocks} \le \#\text{observables}.
\end{equation}
Instead, identification requires local invertibility of the likelihood mapping,
i.e., full column rank of the Jacobian relating structural parameters to
implied second moments.
\end{technicalbox}

Reference: \citet{iskrev2010}.

\section{Post-2008 Financial Friction Models}

\begin{intuitionbox}
Financial-friction models add shocks to match spreads, credit dynamics,
and investment fluctuations. These additions improve empirical fit but increase
the risk of overlapping propagation mechanisms.
\end{intuitionbox}

\begin{technicalbox}
Post-crisis models introduce risk-premium, credit-spread, capital-quality,
and liquidity shocks (e.g., \citealt{christiano2014}; \citealt{gerali2010};
\citealt{gertlerkaradi2011}). Increased shock dimensionality raises the
importance of identification diagnostics and careful interpretation of
variance decompositions.
\end{technicalbox}

\section*{Conclusion}

Having more structural shocks than observables is not inherently problematic.
Identification in DSGE models derives from dynamic propagation and
cross-equation restrictions, not simple counting.
Risks emerge when shocks are observationally equivalent or absorb misspecification.

\bibliographystyle{apalike}
\begin{thebibliography}{}

\bibitem[An and Schorfheide(2007)]{an2007}
An, S. and Schorfheide, F. (2007).
Bayesian analysis of DSGE models.
\textit{Journal of Economic Dynamics and Control}.

\bibitem[Christiano et al.(2014)]{christiano2014}
Christiano, L., Motto, R., and Rostagno, M. (2014).
Risk shocks.
\textit{American Economic Review}.

\bibitem[Cúrdia and Reis(2012)]{curdia2012}
Cúrdia, V. and Reis, R. (2012).
Correlated disturbances and business cycles.
\textit{Journal of Monetary Economics}.

\bibitem[Den Haan and Drechsel(2021)]{denhaan2021asd}
Den Haan, W. and Drechsel, T. (2021).
Agnostic structural disturbances.
\textit{Journal of Monetary Economics}.

\bibitem[Durbin and Koopman(2012)]{durbin2012}
Durbin, J. and Koopman, S. (2012).
\textit{Time Series Analysis by State Space Methods}.
Oxford University Press.

\bibitem[Gerali et al.(2010)]{gerali2010}
Gerali, A. et al. (2010).
Credit and banking in a DSGE model.
\textit{Journal of Money, Credit and Banking}.

\bibitem[Gertler and Karadi(2011)]{gertlerkaradi2011}
Gertler, M. and Karadi, P. (2011).
Unconventional monetary policy.
\textit{Journal of Monetary Economics}.

\bibitem[Hamilton(1994)]{hamilton1994}
Hamilton, J. (1994).
\textit{Time Series Analysis}.
Princeton University Press.

\bibitem[Herbst and Schorfheide(2015)]{herbst2015}
Herbst, E. and Schorfheide, F. (2015).
\textit{Bayesian Estimation of DSGE Models}.
Princeton University Press.

\bibitem[Iskrev(2010)]{iskrev2010}
Iskrev, N. (2010).
Local identification in DSGE models.
\textit{Journal of Monetary Economics}.

\bibitem[Smets and Wouters(2007)]{smets2007}
Smets, F. and Wouters, R. (2007).
Shocks and frictions in US business cycles.
\textit{American Economic Review}.

\end{thebibliography}

\end{document}